\documentclass{ximera}[font=12pt]


\title{Adding Integers: Same Signs}   % Use xmlatex all -s to compile when SERVE refuses to work.
\author{Kelly Stady}
\license{CC: 0}         % replace with an appropriate license, or set it in xmPreamble

\begin{document}

\begin{abstract}

\end{abstract}

\maketitle
\label{xim:AddIntSameSActivity}


\section*{Adding Numbers with the Same Sign}

When adding numbers with the same sign, there is no ``battle.'' Instead, the numbers \textbf{join forces} to create a larger 
asset or a deeper debt.

\begin{example} \textbf{Joining Assets}

If you have \$400 in your savings account ($+400$) and you win \$150 ($+150$), you have a net balance of \$550.
\[ 400 + 150 = 550 \]

\end{example}

\begin{example} \textbf{Joining Debts}

If you owe your brother \$10 ($-10$) and then borrow another \$25 ($-25$), your net balance is a debt of \$35.  So, you are falling deeper into debt.
\[ -10 + (-25) = -35 \]

\end{example}


We can now write a formal mathematical rule for adding integers with the same sign.  Once again, we will use absolute value.


\begin{procedure}[\textbf{Steps for Adding Integers with the Same Signs}]

    ~ \\[2pt]
\begin{enumerate} 
    \item Find the \textbf{absolute value} or \textbf{strength} of each number.

     \item \textbf{Add} the absolute values.

     \item Write the answer with the common sign.
\end{enumerate}
\end{procedure}


\begin{example} Let's use our three-step process to find the sum $-50 + (-30).$

\begin{description}[style=multiline, leftmargin=5em, font=\bfseries]
    \item[Step 1] \textbf{Find the Absolute Values (Strengths):} \\
    Fnd the strength of each number.
    \[ |-50| = 50 \quad \text{and} \quad |30| = 30 \]

    \item[Step 2] \textbf{Add the Absolute Values (Join Forces):} \\
    Since the signs are the same, we add their absolute values.
    \[ 50 + 30 = 80 \]

    \item[Step 3] \textbf{Determine the Sign:} \\
    Since both numbers are negative, the answer is negative.
    \[ -50 + (-30) = -80 \]
\end{description}

\end{example}


\begin{problem}
    Apply the steps to find the sum $-15 + (-40).$

    ~ \\[3em]
    \textbf{Step 1:} Find the \textbf{absolute values} (\textbf{Strengths}):
    \[ |-15| = \answer{15} \quad \text{and} \quad |-40| = \answer{40} \]

    \begin{problem}
    \textbf{Step 2:} \textbf{Add} the absolute values (\textbf{Join Forces}):
    \[ 15 + 40 = \answer{55} \]

    \begin{problem}
    \textbf{Step 3:} \textbf{Determine the sign}:

    Write the answer with the common sign.
    \textbf{Final Answer:} $-15 + (-40) = \answer{-55}$

    \begin{feedback}[correct]
        Good job! You found the strengths, helped the numbers to join forces and wrote the answer with the common sign.
    \end{feedback}

\end{problem}
\end{problem}
\end{problem}



\begin{problem}
Follow the steps below to find the sum.

\[ -80 + (2) \]

Will the numbers join forces or battle?

\begin{multipleChoice}
\choice{Join forces.}
\choice[correct]{Battle.}
\end{multipleChoice}

\begin{problem}
Find the absolute value (strength) of each number.

\[ |-80| \begin{prompt} = \answer{80} \end{prompt} \hspace{1in} |2| \begin{prompt} = \answer{2} \end{prompt} \]

\begin{problem}
How will you find the answer?

\begin{multipleChoice}
\choice{Add the absolute values and write the answer with the common sign.}
\choice{Add the absolute values and write the answer with the sign of the stronger number.}
\choice{Subtract the absolute values and write the answer with the common sign.}
\choice[correct]{Subtract the absolute values and write the answer with the sign of the stronger number.}
\end{multipleChoice}

\begin{problem}
Which number is stronger (has larger absolute value)?

\begin{multipleChoice}
\choice{$2$}
\choice[correct]{$-80$}
\end{multipleChoice}

\textbf{Answer:}  $-80 + 2 = \answer{-78}$

\end{problem}
\end{problem}
\end{problem}
\end{problem}


\begin{problem}
Follow the steps below to find the sum.

\[ -900 + (-75) \]

Will the numbers join forces or battle?

\begin{multipleChoice}
\choice{Battle.}
\choice[correct]{Join forces.}
\end{multipleChoice}

Find the absolute value (strength) of each number.

\[ |-900| \begin{prompt} = \answer{900} \end{prompt} \hspace{1in} |-75| \begin{prompt} = \answer{75} \end{prompt} \]

\begin{problem}
How will you find the answer?

\begin{multipleChoice}
\choice{Add the absolute values and write the answer with the sign of the stronger number.}
\choice[correct]{Add the absolute values and write the answer with the common sign.}
\choice{Subtract the absolute values and write the answer with the sign of the stronger number.}
\choice{Subtract the absolute values and write the answer with the common sign.}
\end{multipleChoice}
 
\begin{problem}


\textbf{Answer:}  $-900 + (-75) = \answer{-975}$

\end{problem}
\end{problem}
\end{problem}


ADD MORE PROBLEMS and VIDEOS

\end{document}
